\begin{notechapter}
\chapter{Generative and Discriminative Models}\label{ch:generative-discriminative}
\chaptersubtitle{Two products of probabilities: likelihood and naive Bayes}

\begin{chapterintro}
Products of probabilities are everywhere in machine learning, and two of them
are easily confused: the product over \emph{data points} in a likelihood, and
the product over \emph{features} in the naive Bayes model. They rest on
different assumptions and play different roles. Disentangling them leads to
the distinction between generative and discriminative models. You will learn
\begin{itemize}[nosep]
  \item where the product in a likelihood comes from, and why maximizing a
        conditional likelihood is the same as minimizing the cross-entropy;
  \item the naive Bayes model, its conditional-independence assumption, and
        how to train it by counting (with a derivation);
  \item a complete worked example of a spam filter;
  \item the difference between the joint and the conditional likelihood;
  \item why Gaussian naive Bayes and logistic regression have the same form
        but are not the same model.
\end{itemize}
\prerequisites{\cref{ch:learning-problem,ch:bayes-decision-theory};
maximum likelihood; Lagrange multipliers (for one proof).}
\end{chapterintro}

\section{Two products}\label{generative-discriminative:sec:two}

Consider a dataset $d_n = \{(x_i,y_i)\}_{i=1}^n$ with inputs
$x_i = (x_{i1},\dots,x_{iD})$ made of $D$ features, and labels in a finite set
$\Ys$. Two expressions appear constantly:
\[
  \underbrace{\prod_{i=1}^n Q_\theta(y_i\mid x_i)}_{\text{a likelihood: product over data points}}
  \qquad\text{and}\qquad
  \underbrace{P(Y = y)\prod_{j=1}^D P(X_j = x_j\mid Y = y)}_{\text{naive Bayes: product over features}} .
\]
The first is a \emph{training principle}; it applies to any probabilistic
model. The second is a \emph{model}: a particular assumption about how the
features of one example relate to each other.
\Cref{generative-discriminative:tab:compare} summarizes the contrast; the rest
of the chapter explains it.

\begin{table}[ht]
\centering
\caption{The two product structures.}
\label{generative-discriminative:tab:compare}
\begin{tabular}{@{}p{0.2\textwidth}p{0.36\textwidth}p{0.36\textwidth}@{}}
\toprule
 & Likelihood product & Naive Bayes product \\
\midrule
multiplies over & data points $i = 1,\dots,n$ & features $j = 1,\dots,D$ of one point \\
comes from & independence of the examples (i.i.d.) & conditional independence of features given the class \\
role & how to fit any model & what one particular model assumes \\
status & standard, usually reasonable & strong, usually false, often harmless \\
\bottomrule
\end{tabular}
\end{table}

\section{The likelihood: a product over data points}\label{generative-discriminative:sec:likelihood}

Let $Q_\theta(y\mid x)$ be a model of the conditional distribution of $Y$ given
$X$, with parameters $\theta$ --- for instance logistic regression or a neural
network with a softmax output.

\begin{definition}[Conditional likelihood]
The \emph{conditional likelihood} of $\theta$ on the dataset is the probability
the model assigns to the observed labels given the observed inputs,
\[
  \mathcal L(\theta) = \prod_{i=1}^n Q_\theta(y_i\mid x_i),
  \qquad
  \log\mathcal L(\theta) = \sum_{i=1}^n\log Q_\theta(y_i\mid x_i).
\]
\end{definition}

The product is the probability of the $n$ labels \emph{jointly}, and it
factorizes because the examples are assumed independent: by the i.i.d.\
assumption of \cref{learning-problem:def:sample},
$P(y_1,\dots,y_n\mid x_1,\dots,x_n) = \prod_i P(y_i\mid x_i)$. Nothing is
assumed about the features inside each $x_i$.

Taking logarithms turns the product into a sum, which is numerically stable
and easier to differentiate. Dividing by $-n$ gives
\[
  -\frac1n\log\mathcal L(\theta) = \frac1n\sum_{i=1}^n\bigl(-\log Q_\theta(y_i\mid x_i)\bigr),
\]
which is the empirical risk of $Q_\theta$ under the log-loss of
\cref{learning-problem:tab:losses}, also called the \emph{cross-entropy}.
\textbf{Maximizing the conditional likelihood is ERM with the log-loss}, and by
\cref{bayes-decision-theory:prop:logloss} its population version is minimized
by the true conditional distribution. This is why the same loss is used for
logistic regression and for deep classifiers.

\section{Naive Bayes: a product over features}\label{generative-discriminative:sec:nb}

Naive Bayes is a \emph{generative} classifier: instead of modelling
$P(Y\mid X)$, it models how each class generates its inputs and inverts the
model with Bayes' rule.

\begin{definition}[Naive Bayes]
The naive Bayes model assumes that the features are \emph{conditionally
independent given the class}:
\[
  P(X_1 = x_1,\dots,X_D = x_D\mid Y = y) = \prod_{j=1}^D P(X_j = x_j\mid Y = y).
\]
Combined with Bayes' rule, the posterior is
\[
  P(Y = y\mid X = x) = \frac{P(Y=y)\prod_j P(X_j = x_j\mid Y = y)}{\sum_{y'}P(Y=y')\prod_jP(X_j = x_j\mid Y = y')}
  \;\propto\; P(Y = y)\prod_{j=1}^D P(X_j = x_j\mid Y = y),
\]
and the classifier predicts the class with the largest posterior (the MAP rule of
\cref{bayes-decision-theory:sec:map-ml}).
\end{definition}

\begin{intuition}[Why ``naive''?]
In a spam filter the words ``money'' and ``free'' are not independent, even
within spam: they tend to appear together. Naive Bayes ignores this, and
counts their evidence twice. The probabilities it outputs are therefore often
too extreme. But the classification only needs the \emph{ranking} of the
classes to be right, and that is far more robust: naive Bayes is frequently a
surprisingly good classifier with badly calibrated probabilities.
\end{intuition}

The benefit of the assumption is statistical. A full joint model of $D$ binary
features needs $2^D - 1$ parameters per class; naive Bayes needs $D$. That is
the difference between an impossible estimation problem and a trivial one.

\subsection{Training naive Bayes by maximum likelihood}

A generative model specifies the whole joint distribution
$P_\theta(x,y) = P_\theta(y)P_\theta(x\mid y)$, and is naturally fitted by
maximizing the \emph{joint} likelihood. For naive Bayes, the two products now
appear together:
\[
  \mathcal L_{\mathrm{joint}}(\theta) = \prod_{i=1}^n P_\theta(x_i, y_i)
  = \prod_{i=1}^n\Bigl[P_\theta(y_i)\prod_{j=1}^D P_\theta(x_{ij}\mid y_i)\Bigr].
\]
The outer product comes from the independence of the examples, the inner one
from the model.

\begin{proposition}[Maximum-likelihood estimates are frequencies]\label{generative-discriminative:prop:counts}
For discrete features, the joint likelihood is maximized by
\[
  \hat P(Y = y) = \frac{N_y}{n},\qquad
  \hat P(X_j = v\mid Y = y) = \frac{N_{j,v,y}}{N_y},
\]
where $N_y$ is the number of examples of class $y$ and $N_{j,v,y}$ the number of
examples of class $y$ whose feature $j$ equals $v$.
\end{proposition}

\begin{proof}
The log-likelihood is
\[
  \log\mathcal L_{\mathrm{joint}} = \sum_{y} N_y\log\pi_y
  + \sum_{j}\sum_{y}\sum_v N_{j,v,y}\log\phi_{j,v,y},
\]
with $\pi_y = P(Y = y)$ and $\phi_{j,v,y} = P(X_j = v\mid Y = y)$: each example
contributes the log-probability of its class and of each of its feature values.
The parameters separate into independent groups --- the priors $(\pi_y)_y$, and
for each pair $(j,y)$ the distribution $(\phi_{j,v,y})_v$ --- each constrained
to sum to $1$. Each group is a problem of the form
\[
  \max_{p}\ \sum_v c_v\log p_v\quad\text{subject to}\quad \sum_v p_v = 1,
\]
with counts $c_v\geq0$. Setting the derivative of the Lagrangian
$\sum_v c_v\log p_v - \lambda(\sum_v p_v - 1)$ to zero gives $c_v/p_v = \lambda$,
so $p_v\propto c_v$, and the constraint gives $p_v = c_v/\sum_{v'}c_{v'}$. The
objective is concave, so this is the maximum.
\end{proof}

\begin{pitfall}[Zero counts]
If the word ``refund'' never occurs in the training ham, the estimate
$\hat P(\text{refund}\mid\text{ham}) = 0$ makes every future email containing
``refund'' certain spam, whatever else it says. The standard fix is
\emph{Laplace} (add-one) smoothing:
$\hat P(X_j = v\mid Y = y) = (N_{j,v,y} + 1)/(N_y + V_j)$, where $V_j$ is the
number of values of feature $j$. It corresponds to a Bayesian estimate with a
uniform prior on the parameters (\cref{ch:bayesian-reasoning}).
\end{pitfall}

\section{A worked example: a tiny spam filter}\label{generative-discriminative:sec:spam}

\begin{example}[Two words, ten emails]\label{generative-discriminative:ex:spam}
Each email is described by two binary features: $M = 1$ if it contains
``money'' and $F = 1$ if it contains ``free''. The training set contains
four spam and six ham emails:
\begin{center}
\begin{tabular}{@{}lcccc@{}}
\toprule
class & $(M,F)$ of each email & $N_y$ & $\#\{M=1\}$ & $\#\{F=1\}$ \\
\midrule
spam & $(1,1)$, $(1,0)$, $(0,1)$, $(1,1)$ & 4 & 3 & 3 \\
ham  & $(0,0)$, $(1,0)$, $(0,0)$, $(0,1)$, $(0,0)$, $(0,0)$ & 6 & 1 & 1 \\
\bottomrule
\end{tabular}
\end{center}
\emph{Training} (\cref{generative-discriminative:prop:counts}):
\[
  \hat P(\text{spam}) = \tfrac{4}{10},\quad
  \hat P(M{=}1\mid\text{spam}) = \hat P(F{=}1\mid\text{spam}) = \tfrac34,\quad
  \hat P(M{=}1\mid\text{ham}) = \hat P(F{=}1\mid\text{ham}) = \tfrac16 .
\]
\emph{Prediction} for a new email containing both words:
\begin{align*}
  \text{spam}:&\quad 0.4\times\tfrac34\times\tfrac34 = 0.225,\\
  \text{ham}:&\quad 0.6\times\tfrac16\times\tfrac16 \approx 0.0167,
\end{align*}
so $\hat P(\text{spam}\mid M{=}1,F{=}1) = 0.225/(0.225 + 0.0167)\approx 0.93$.
For an email containing neither word, the scores are
$0.4\times\frac14\times\frac14 = 0.025$ and $0.6\times\frac56\times\frac56\approx0.417$,
giving $\hat P(\text{spam}\mid M{=}0,F{=}0)\approx0.06$.

The \emph{product over features} was used to score one email; the \emph{product
over examples} was used, implicitly, to obtain the counts.
\end{example}

\section{Joint versus conditional likelihood}\label{generative-discriminative:sec:joint-vs-conditional}

The two likelihoods of this chapter are related by
\[
  \prod_i P_\theta(x_i,y_i) = \underbrace{\prod_i P_\theta(y_i\mid x_i)}_{\text{conditional likelihood}}\;\underbrace{\prod_iP_\theta(x_i)}_{\text{likelihood of the inputs}} .
\]
\begin{itemize}
  \item A \textbf{generative} model (naive Bayes, Gaussian discriminant
        analysis, hidden Markov models) maximizes the joint likelihood. It
        spends part of its effort modelling the inputs, which classification
        does not need, but it can generate data, handle missing features, and
        often has closed-form estimates.
  \item A \textbf{discriminative} model (logistic regression, neural networks,
        conditional random fields) maximizes the conditional likelihood only
        --- the cross-entropy. It models exactly what prediction requires.
\end{itemize}
The closed-form counts of \cref{generative-discriminative:prop:counts} maximize
the joint likelihood. Maximizing instead the conditional likelihood of the
\emph{same} naive Bayes model has no closed form --- and, as the next section
shows, amounts to logistic regression.

\begin{figure}[ht]
\centering
\begin{tikzpicture}[font=\small, node distance=8mm and 4mm,
  box/.style={draw, rounded corners, align=center, inner sep=4pt, fill=#1!6, draw=#1},
  >=Stealth]
  \node[box=bookaccent] (mle) {\textbf{maximum likelihood}\\ (product over examples)};
  \node[box=bookdef, below left=10mm and 4mm of mle] (gen) {\textbf{generative}: joint likelihood\\ $\max\prod_iP_\theta(x_i,y_i)$\\[2pt] naive Bayes, GDA, HMM};
  \node[box=bookexample, below right=10mm and 4mm of mle] (disc) {\textbf{discriminative}: conditional likelihood\\ $\max\prod_iQ_\theta(y_i\mid x_i)$ = min cross-entropy\\[2pt] logistic regression, neural nets, CRF};
  \draw[->] (mle) -- (gen);
  \draw[->] (mle) -- (disc);
\end{tikzpicture}
\caption{Maximum likelihood is the common training principle; generative and
discriminative models differ in \emph{which} likelihood they maximize.}
\label{generative-discriminative:fig:hierarchy}
\end{figure}

\section{Gaussian naive Bayes and logistic regression}\label{generative-discriminative:sec:gnb-lr}

Take two classes and real features, and model each feature, given the class,
as a Gaussian with a class-dependent mean and a variance \emph{shared by the
two classes}:
$X_j\mid Y = c\sim\mathcal N(\mu_{jc},\sigma_j^2)$. The log-odds are
\[
  \log\frac{P(Y=1\mid x)}{P(Y=0\mid x)}
  = \log\frac{\pi_1}{\pi_0} + \sum_{j=1}^D\log\frac{\mathcal N(x_j;\mu_{j1},\sigma_j^2)}{\mathcal N(x_j;\mu_{j0},\sigma_j^2)} .
\]
In each term the normalizing constants cancel (same variance), leaving
\[
  \frac{(x_j-\mu_{j0})^2 - (x_j-\mu_{j1})^2}{2\sigma_j^2}
  = \frac{\mu_{j1}-\mu_{j0}}{\sigma_j^2}\,x_j + \frac{\mu_{j0}^2-\mu_{j1}^2}{2\sigma_j^2}:
\]
the quadratic terms $x_j^2$ cancel. Hence
\[
  \log\frac{P(Y=1\mid x)}{P(Y=0\mid x)} = w_0 + \sum_{j=1}^Dw_jx_j,
  \qquad w_j = \frac{\mu_{j1}-\mu_{j0}}{\sigma_j^2},\quad
  w_0 = \log\frac{\pi_1}{\pi_0} + \sum_j\frac{\mu_{j0}^2-\mu_{j1}^2}{2\sigma_j^2},
\]
that is, $P(Y=1\mid x) = \sigma(w_0 + w^\top x)$ --- exactly the form of logistic
regression. With class-dependent variances the $x_j^2$ terms survive and the
log-odds become quadratic.

\begin{keyidea}[Same form, different fits]
Gaussian naive Bayes and logistic regression describe the same family of
posteriors, but choose a member of it differently. Naive Bayes estimates means,
variances and priors from the joint likelihood and \emph{derives} $w$;
logistic regression fits $w$ \emph{directly} by maximizing the conditional
likelihood. When the naive Bayes assumptions are wrong, logistic regression
is asymptotically better, because it does not rely on them. With little data,
naive Bayes often wins, because its estimates converge faster (Ng and Jordan,
2002). They form a \emph{generative--discriminative pair}.
\end{keyidea}

\section{Choosing between the two approaches}\label{generative-discriminative:sec:choosing}

Naive Bayes (and generative models generally) are attractive when
\begin{itemize}[nosep]
  \item the dataset is small, or training must be very fast;
  \item features are roughly conditionally independent (bag-of-words text is
        the classic case);
  \item a simple, interpretable baseline is wanted, or features may be
        missing at prediction time.
\end{itemize}
Discriminative models trained with the cross-entropy are preferable when
\begin{itemize}[nosep]
  \item data are plentiful and features interact in complex ways;
  \item the best predictive accuracy is the goal;
  \item calibrated probabilities $P(Y\mid X)$ are needed (for instance to apply
        a cost-sensitive threshold, \cref{bayes-decision-theory:thm:costs}).
\end{itemize}

\begin{chaptersummary}
\begin{itemize}[nosep]
  \item The likelihood $\prod_iQ_\theta(y_i\mid x_i)$ is a product over
        \emph{examples}, justified by the i.i.d.\ assumption; maximizing it is
        minimizing the cross-entropy, i.e.\ ERM with the log-loss.
  \item Naive Bayes' $P(y)\prod_jP(x_j\mid y)$ is a product over
        \emph{features}, justified only by the model's conditional-independence
        assumption.
  \item Training naive Bayes by maximum (joint) likelihood gives frequency
        estimates; Laplace smoothing avoids zero probabilities.
  \item Generative models maximize the joint likelihood, discriminative ones
        the conditional likelihood. Gaussian naive Bayes with shared variances
        and logistic regression have the same posterior form but different
        estimators.
\end{itemize}
\end{chaptersummary}

\section*{Exercises}
\addcontentsline{toc}{section}{Exercises}

\begin{exercise}
Redo \cref{generative-discriminative:ex:spam} with Laplace smoothing. How do the
two posterior probabilities change?
\end{exercise}

\begin{exercise}
Show that for binary features, the log-odds of a naive Bayes model are also an
affine function of $x$. Express the weights in terms of
$\phi_{j1} = P(X_j = 1\mid Y=1)$ and $\phi_{j0} = P(X_j=1\mid Y=0)$.
\hint{write $P(X_j=x_j\mid Y=c) = \phi_{jc}^{x_j}(1-\phi_{jc})^{1-x_j}$.}
\end{exercise}

\begin{exercise}
Duplicate the feature ``money'' in \cref{generative-discriminative:ex:spam}
(so that each email has features $(M, M, F)$). Recompute the posterior for an
email with $M = 1$, $F = 0$ before and after duplication, and explain the effect
in terms of the conditional-independence assumption.
\end{exercise}
\end{notechapter}
